\section{Carrying window updates through cached potential modes}
\label{sec:window-modes}

Matching updates use only $d_0+2$ parameters, but the preceding version
still projects the full source potential matrix onto each updated wide
Q basis. This continuation carries the same thin row operations through
cached source modes. It computes exactly the same corner forms, then uses
the existing Euler screen, standard geometry and independent disc observers.

\subsection{Corrected columns and augmented potential modes}

Retain the decompositions $A_j=A_Sc_j+r_j$ of
Section~\ref{sec:window-basis}, with the source potential matrix $P$.
For each nonbase type define its algebraic correction
\[
                 D_j=P_j-P_Sc_j.
\]
The old matching basis contributes $P_SK_0^T$. A zero-residual singleton
adds $D_j$; two such columns add their separate corrections. An opposite
pair adds $D_a+\gamma D_b$, with the same positive ratio satisfying
$r_a+\gamma r_b=0$ as in the matching update.
These are precisely the potentials of the augmented matching modes.
A correction with nonzero residual is not a matching vector by itself;
it becomes part of the actual matching space only through the cancelling
pair and subsequent restrictions. No individual correction is asserted
to be an admissible normal surface.

Write the augmented matching rows in full global Q coordinates as $W$,
and their potential columns as $\Phi_W=PW^T$. Replacement constraints give
a thin row matrix $N$. Sorting the surviving support changes coordinate
order but not the full global vector. Canonical gauge reduction gives a
thin matrix $G$. Thus, before suppressing zero removed coordinates,
\[
              K'=GNW,\qquad
              P(K')^T=\Phi_WN^TG^T.
\]
The implementation retains $Z=GN$, then combines the cached columns of
$\Phi_W$ with its rows.

\begin{proposition}[Exact propagated source potentials]
The propagated corner forms equal direct multiplication of the source
potentials by the updated canonical Q basis, at every physical corner.
Consequently their equal-form identifications, canonical lifts, Euler
screens and complete standard matching cones agree exactly.
\end{proposition}
\begin{proof}
For every augmented mode its potential is the corresponding column of
$PW^T$, by the linearity of $P$ and the displayed corrected-column formulas.
Replacement and gauge reduction are ordinary linear combinations of these
modes, so multiplication commutes with both operations. Removed coordinates
are exactly zero after replacement constraints; dropping them therefore
changes neither the full global Q vector nor its potential. The displayed
matrix equality proves the corner-by-corner identity. Subsequent
identification, minimum normalization and matching construction use only
these exact forms and the unchanged source, giving the remaining claims.
\end{proof}

\subsection{Track the actual gauge operations}

During the native reversed-column gauge reduction, append an identity
matrix to the thin input rows. Only the original matching columns are
eligible pivots. Every elimination operation also acts on the appended
columns, recording the exact coefficients of the resulting rows in the
input mode basis. Reversing the matching columns back and sorting by free
index preserves the corresponding transformation rows. Composing this
record with the replacement nullspace gives $Z$.

Auxiliary identity columns never become matching variables or pivot choices.
They do not change the canonical gauge. The untracked entry point continues
to return the same ordinary support/basis pair. The window producer requests
an additional private projection plan, which remains outside every portable
positive certificate.

\subsection{Lazy sparse source caches}

The old basis potentials are projected at most once and stored by nonzero
corner entries. The all-Q potential rows are transposed at most once into
sparse columns. Each needed $D_j$ is then accumulated from $P_j$ and the
nonzero coefficients of $c_j$, with exact zero cancellations removed.
A correction is cached only after its complete construction. An interrupted
transpose or column cannot publish an unfinished vector for another query.

Per-window projection combines only the nonzero thin weights and sparse
cached potential entries. It reconstructs a full tuple of corner forms for
the existing Euler and standard-kernel interfaces. Algebraic zeros are
represented exactly; no numerical threshold or hashing surrogate changes
form equality. Source caches belong to one read-only source and fixed old
support. They are not shared across unrelated centres or triangulations.

Projection remains demand driven. An empty updated kernel or an immediately
empty one-dimensional nonnegative cone returns before creating a cache.
Unneeded correction columns are never computed. A positive result and all
independent source replay retain their previous contract. High-nullity
arrangement and support fallbacks use the identical propagated forms.

\subsection{Arithmetic cost and retained scope}

Let $C$ be the number of distinct correction columns actually needed and
$m\le d_0+2$ the augmented parameter count. Transposition costs
$O(\operatorname{nnz}P)$ once. The old basis projection costs
$O(\operatorname{nnz}P+d_0\operatorname{nnz}P_S)$ once. A correction's sparse
arithmetic costs
\[
 O\!\left(1+\|c_j\|_0+\operatorname{nnz}P_j+
       \sum_{i:(c_j)_i\ne0}\operatorname{nnz}P_{S_i}\right),
\]
in addition to constructing its coefficient record. Its worst case is
$O(T(k+1))$, while $C\le3T$. The complete correction preparation is
therefore polynomial, bounded by $O(CT(k+1))$ arithmetic operations.

After these preparations, each window uses thin matching/gauge operations
and sparse combinations of at most $d_0+2$ source potential vectors per
output basis row. Including dense interface reconstruction, potential
propagation costs $O(Td(d_0+2))$ in the worst case. It removes repeated
scanning of $\operatorname{nnz}P$ and repeated wide source projection from
each candidate. The lazy cache can be more expensive on a short query that
never reuses a column; the measured tradeoff must include its preparation.

Cached coefficient bits, sparse-vector storage, exact rational arithmetic,
output and independent replay remain charged. This changes a polynomial
factor in the existing local logarithmic-nullity/radius bound. It does not
supply global sector centres, a universal distance bound or a general
quasi-polynomial recognition theorem. Shared deterministic guard counts may
change separately from wall-clock time.

\subsection{Validation and measured behavior}
All 1,408 maintained tests pass in 241.501 seconds; the twelve focused
tests pass in 19.293 seconds. Every tested small literal source window
compares propagated forms with direct source projection, full bases with
fresh reconstruction, and complete native ray sets and coordinate lifts.
A genuine nullity-four source retains both complete generic methods.
An interrupted column construction leaves no partial cache entry, and a
real diagram window still produces an independently replayed ordinary disc.

The 85-source audit exactly preserves all disabled outputs, all enabled
verdicts and all 31 native positives. Seven window discs receive fresh
Regina confirmation. The twenty-eight work-capped cases and twenty-six
completed window misses remain unchanged. On the genus-one source, thirteen
actual projections use one old-basis preparation, one transpose and twenty-four
corrected columns. The trefoil and figure-eight use thirty and forty-two
projections with fifty-four and seventy-six corrected columns, respectively.
Their candidate and emitted-ray counts remain unchanged. Source guard calls
increase from 109,610 to 110,981 on the genus-one case, from 178,409 to 180,453
on the trefoil, and from 314,978 to 318,797 on the figure-eight. Cache preparation
is genuine work; the improved arithmetic organization is not a reduction
in every guard-count allowance. All 599 captured runtime pins stay fixed.
The 249.386-second audit is not a paired speedup estimate.

The first dense-mode pilot completes 120 measured calls and forty warmups
and is mostly near parity. It exposes the cost of multiplying algebraic zero
potential entries. The final version stores sparse modes and accumulates
only their nonzero contributions before reconstructing the existing dense
corner-form interface. The pilot records and both exact historical source
snapshots needed to recover all 599 pins are retained.

The final native comparison freezes four preceding-release modules at
\code{8a6cc03d7}: window basis construction, residual search, seed queries
and recognition. Local imports are redirected only to corresponding frozen
dependencies. Five serial interleaved four-arm rounds complete 200 measured
calls and forty warmups. Native timings include source/cache preparation,
all configured queries, independent positive replay and serialized output.
Both arms finish all configured work and agree on final statuses and tested
positive certificate digests.

\begin{samepage}
\noindent\textbf{Complete native queries with radius-two windows.}
\begin{center}
\small
\begin{tabular}{lrrrrr}
\toprule
Input & previous ms & current ms & old/new & old A/A & new A/A\\
\midrule
empty circle & 6.039 & 5.995 & 1.012 & 1.004 & 1.005\\
optimized positive & 59.252 & 100.487 & 0.927 & 0.965 & 0.989\\
genus one miss & 240.396 & 234.586 & 1.016 & 1.005 & 0.989\\
trefoil & 408.116 & 403.060 & 1.015 & 0.982 & 1.006\\
figure eight & 751.056 & 722.209 & 1.038 & 1.010 & 0.983\\
\bottomrule
\end{tabular}
\end{center}

\end{samepage}

The optimized-positive native control returns before mode projection, yet its
first comparison has a paired old/new ratio of 0.927 and noticeably different
raw medians. This is not attributed to the unused cache implementation.
A fifteen-round repeat completes sixty control calls with four warmups,
a ratio of 0.993 and A/A controls 0.989 and 1.003, resolving that row near
parity. Its source evidence confirms that neither arm enters window search.

\begin{samepage}
\noindent\textbf{Complete recognition with radius-two windows enabled.}
\begin{center}
\small
\begin{tabular}{lrrrrr}
\toprule
Input & previous ms & current ms & old/new & old A/A & new A/A\\
\midrule
empty circle & 0.018 & 0.018 & 1.011 & 1.000 & 0.952\\
optimized positive & 51.464 & 51.683 & 0.995 & 0.989 & 1.011\\
genus one miss & 246.124 & 238.624 & 1.015 & 1.000 & 0.983\\
trefoil & 444.265 & 428.993 & 1.038 & 1.013 & 1.007\\
figure eight & 777.020 & 763.385 & 1.019 & 0.983 & 1.004\\
\bottomrule
\end{tabular}
\end{center}

\end{samepage}

The three window cases have paired complete-recognition ratios of 1.015,
1.038 and 1.019 on the genus-one, trefoil and figure-eight, respectively.
These are modest small-case measurements. The same longer repeat performs
sixty complete trefoil recognition calls and four warmups, including the
established complete fallback after the native miss. Its gain is 1.015,
with A/A controls 0.998 and 1.002, smaller than the first row's gain.
Both runs remain available. Early full-recognition controls stay near parity;
the microsecond empty case is not used to infer a projection speedup.
The principal result is the exact amortized projection bound, not a uniform
wall-clock gain. Radius zero remains much faster on these small controls,
so the window stage continues to default off.

\begin{sloppypar}
Runtime release \code{4f181d86c}, all tests, the source audit, dense pilot,
final timing run and 120-call control/recognition repeat are retained.
From \code{fast/}, use \code{normal\_orbit\_research.window\_modes} with
\code{audit --fresh-regina} or \code{benchmark --rounds 5}.
Publication checks verify 599 frozen/current runtime pins, 598 preceding
release evidence pins, the recoverable pilot and every final timing outcome.
They independently replay all 31 saved positive source proofs with mode
projection, corrected-column caches, Euler compilation, prepared observers,
matching updates, residual planning, coherent extraction, ray enumeration
and orbit producers disabled. Independent source reconstruction remains
the authority for positive results.
\end{sloppypar}

